\documentclass[a4paper,11pt,openany]{book}
\usepackage[utf8]{inputenc}
\usepackage[left=2.54cm,top=2.54cm,right=2.54cm,bottom=2.54cm]{geometry}
\usepackage[spanish]{babel}
\usepackage{amsmath}
\usepackage{amsfonts}
\usepackage{amssymb}
\usepackage{graphicx}
\usepackage{color}
\usepackage[usenames,dvipsnames]{xcolor}
\usepackage{pifont}
\usepackage{marvosym}
\newtheorem{teo}{Teorema}
\newtheorem{ejemplo}{Ejemplo}
\newtheorem{defi}{Definición}
\newtheorem{coro}{Corolario}
\newtheorem{prueba}{Prueba}
\newtheorem{exmp}{Example}[section]
\newtheorem{ejer}{Ejercicio}[section]
\def\proof{\paragraph{\textsf{Demostración.} }}
\def\endproof{\hfill $\blacksquare$ \\}
\usepackage{multirow, array} % para las tablas
\usepackage{multirow}
\usepackage{tabularx}
\usepackage{float} % para usar [H]
\usepackage{tikz}
\usepackage[all]{xy}
\usepackage{cancel}
\usetikzlibrary{positioning}
\usepackage{enumitem}
\newcommand*{\itembolasazules}[1]{% bolas 3D
\footnotesize\protect\tikz[baseline=-3pt]%
\protect\node[scale=.7, circle, shade, ball
color=green]{\color{white}\Large\bf#1};}
\usepackage{tcolorbox} 
\tcbuselibrary{listingsutf8}
\newtcolorbox[auto counter,number within=section]{example}[2][]
{colback=green!5!white,colframe=green!75!black,fonttitle=\bfseries, title=Ejercicio~\thetcbcounter: #2,#1}
\usepackage{background}
\backgroundsetup{
placement=center,
angle=0,
scale=1.1,
contents= {{\includegraphics{HojaCuadriculada.png}}}
}
 
 
\begin{document}
\begin{titlepage}
 
\begin{center}
\vspace*{-1in}
\begin{figure}[htb]
\begin{center}
\includegraphics[width=7cm]{ETITC.png}
\end{center}
\end{figure}
 
 
{\sc \huge Escuela Tecnológica Instituto Técnico Central (ETITC)}\\
\vspace*{0.15in}
Facultad de sistemas\\
\vspace*{0.6in}
\begin{Large}
\textbf{Taller 1: Operaciones Básicas con Números Complejos} \\
\textbf{Matem{\'a}ticas Especiales}\\
\end{Large}
\vspace*{0.3in}
\begin{large}
{\bf Autores} \\
 
\ 
 
Sergio Alejandro Enrrique Caballero Leon\\ 
Johan Alejandro Sogamoso Camacho \\ 
\end{large}
\vspace*{0.3in}
 
\end{center}
 
\begin{center}
{\bf Presentado a:} \\
 
\ 
 
Carlos Romero \\
 
\
 
Bogot{\'a}, 21 de Agosto de 2022.
\end{center}
 
\end{titlepage}

\newpage

\definecolor{ao(english)}{rgb}{0.0, 0.5, 0.0}

\begin{center}
\textbf{Gráfica y Módulo}
\end{center}

Dado los números $\bf{z_{1}}$ y $\bf{z_{2}}$, en cada uno de los ejercicios \ding{172} al \ding{174} grafique en un plano complejo los números $\bf{z_{1}}$, $\bf{z_{2}}$, $\bf{z_{1}\,+\,z_{2}}$ y $\bf{z_{1}\,-\,z_{2}}$. Finalmente calcule el módulo de los cuatro números.\\

\textcolor{ao(english)}{\ding{172}}\,\quad\,$\bf{z_{1}\,=\,6\,-\,2\,i \quad\,;\,\quad z_{2}\,=\,2\,-\,5\,i}$

\textcolor{ao(english)}{\ding{173}}\,\quad\,$\bf{z_{1}\,=\,-\,5\,+\,5\,i \quad\,;\,\quad z_{2}\,=\,-\,6\,-\,4\,i}$

\textcolor{ao(english)}{\ding{174}}\,\quad\,$\bf{z_{1}\,=\,8\,-\,6\,i \quad\,;\,\quad z_{2}\,=\,-\,4\,-\,8\,i}$

\begin{center}
\textbf{Operaciones Aritméticas}
\end{center}

En cada uno de los ejercicios \ding{175} al \ding{179} realice las operaciones indicadas simplificando tanto como sea posible (Recuerde que: $\bf{i^{2}\,=\,-\,1}$)\\

\textcolor{ao(english)}{\ding{175}}\,\quad\,$\bf{3\left(-\,1\,+\,4\,i\right)\,-\,2\left(7\,-\,i\right)}$

\textcolor{ao(english)}{\ding{176}}\,\quad\,$\bf{\left(i\,-\,2\right)\,\left[2\left(1\,+\,i\right)\,-\,3\left(i\,-\,1\right)\right]}$

$$$$

\textcolor{ao(english)}{\ding{177}}\,\quad\,$\bf{\dfrac{\left(2\,+\,i\right)\,\left(3\,-\,2\,i\right)\,\left(1\,+\,2\,i\right)}{\left(1\,-\,i\right)^{2}}}$

$$$$

\textcolor{ao(english)}{\ding{178}}\,\quad\,$\bf{\left(2\,i\,-\,1\right)^{2}\,\left[\dfrac{4}{1\,-\,i}\,+\,\dfrac{2\,-\,i}{1\,+\,i}\right]}$

$$$$

\textcolor{ao(english)}{\ding{179}}\,\quad\,$\bf{\dfrac{i^{4}\,+\,i^{9}\,+\,i^{16}}{2\,-\,i^{5}\,+\,i^{10}\,-\,i^{15}}}$

\end{document}